% s-206-p206-m003; original PDF 206, 207, 208
\Aphorism{s-206-p206-m003}{299}
Der Begriff des einmaligen unumkehrbaren Geschehens ist in unsere Ästhetik des Tragischen einzuführen: er ist ihr sehr nöthig. Nicht nur unser Wissen muss eine Wahrheit, sondern unser Empfinden soll eine Gerechtigkeit darin sehen, dass eine Flamme brennt und zehrt, dass eine That etwas kostet (den Thäter!), dass jedes Glück seinen Gipfel und sein Ende, jedes Leben den Tod sucht. Diese Thatbestände und Nothwendigkeiten gelten im bürgerlichen Leben für leidvoll; aber daraus zu schliessen, dass sie in der Tragödie dem \Sperr{Mitleiden} vorgeführt würden — dieser Irrthum dürfte nur einer Ästhetik begegnen, die sich am Schauspiel eines Feuerwerks ohne Verbrennung, einer Activität ohne Hingabe, eines Lebens ohne Wahn und Opfer, eines Glückes ohne Höhenpunkte und Entwicklung erfreuen kann. Der Normalmensch begreift ein folgerichtiges, wohlgezieltes, rücksichtsloses Geschehen nicht, weil er es im eignen Leben nicht wiederfindet: ihm, der sich zwischen dem Ja des jetzigen und dem Nein des nächsten Augenblicks hundert Mittelwege und Hinterthüren offen lässt, muss die stolze Bereitwilligkeit, mit der in der Tragödie gegeben und genommen, gelebt und gestorben, Ziele erstrebt und Folgen getragen werden, unnöthig hart und grausam erscheinen. Er sieht Leiden und spendet Mitleid, wo für den heroischen Menschen die eigentliche Lebensfreude beginnt oder wo zum mindesten eine Nothwendigkeit waltet: — dieselbe Nothwendigkeit, grob bildlich gesprochen, die für jeden ehrlichen Käufer vorliegt, die Nothwendigkeit zu zahlen, zu opfern, das Eine um des Anderen willen dahinzugeben. Man kann nicht zugleich die Waare und das Geld behalten, nicht zugleich das Ziel treffen und den Pfeil sparen: man kann auch nicht zugleich leben und die Ehre, oder Recht, oder Ruhm behalten, zugleich leben und von der Schuld seines Lebens entsühnt werden, zugleich leben und einer Offenbarung ansichtig werden, zugleich leben und in seiner Liebe selig sein. Das sind Schicksale, die man verstehen, aber nicht bemitleiden oder fürchten soll, erhabene Bilder und Zeichen am Lebenswege, Nothwendigkeiten, deren Erkenntniss für den verwandten Menschen schmerzlos ist: — überdies rein innere Nothwendigkeiten, die keines äusseren Angriffspunktes bedürfen. Das Tragische kommt aus der Sphäre des Individuums selbst und würde zu seinem Untergange führen, wenn es auch gar keine Umwelt gäbe; der Mensch, den ein Ziegelstein, ein Zufall, der Pöbel, der Geist seiner Zeit oder sonst etwas umbringt, ist bedauerlich, — nicht tragisch.
\par

